%=====================================================================
% 第 2 节  综述
% 原文：Section 2. Review（arXiv:1702.06499v2；含 5 个子节，页码从略）
%=====================================================================
\section{综述}

\begin{keytakeaways}
\begin{itemize}
  \item \textbf{物理图像：}本节梳理了通向几何 Langlands 对应的三条进路。Beilinson--Drinfeld 的进路以共形场论为出发点，把对应实现为临界水平仿射李代数的共形块；Kapustin--Witten 的进路把 $N=4$ SYM 理论放在 $\Sigma\times C$ 上，沿 $C$ 紧化得到以 Hitchin 模空间为目标流形的二维 sigma 模型，并把 S-对偶实现为 SYZ 镜像对称；Nekrasov--Witten 的进路用 $\Om$-形变下的 $\CS$ 类理论，把 Virasoro 共形块实现为典范余迷向膜与 oper 膜之间的开弦态空间。
  \item \textbf{数学工具：}$\CD$-模与局部系（含 opers）、Hitchin 模空间及其环面纤维化、SYZ 镜像对称与膜的范畴、Hecke 函子与 Hecke 本征值性质、Virasoro 共形块与 Verlinde 线算符。
  \item \textbf{关键结论：}三条进路有共同的根源——六维 $(2,0)$ 理论在 $\CM^6=\CM^4\times C$ 上的紧化——但两种约化所用的拓扑扭转不同；这一差别既使结果的比较变得微妙，又恰好是把不同描述联系起来的杠杆（Gaiotto--Witten 对边界条件的分类、Betti 几何 Langlands、Cordova--Jafferis 的复 Chern--Simons 描述）。
  \item \textbf{记号提醒：}本节建立了贯穿全文的记号 $\Bun_G$、${}^{L}G$、$\CM_H(G)$、opers 与 $\Om$-形变参数 $\ep_1,\ep_2$（及 $\hbar=\ep_1/\ep_2$）；注意 Kapustin--Witten 与 Nekrasov--Witten 两处的 $\mathsf{B}_{\rm cc}$ 相似但不相同（见 §2.3 脚注），且镜像对称在交换 $G\leftrightarrow{}^{L}G$ 的同时把 $\hbar$ 换成 $1/\hbar$。
\end{itemize}
\end{keytakeaways}

\bigskip
\noindent\rule{\textwidth}{0.6pt}
\bigskip

\subsection{什么是几何 Langlands 对应？}

几何 Langlands 对应通常被宽泛地表述为如下对应：为 Riemann 面 $C$ 上的 ${}^{L}G$-局部系\glos 指派 $\Bun_G$ 上的 $\CD$-模。这里 ${}^{L}G$ 是单复李群 $G$ 的 Langlands 对偶群。这一对应中出现的 ${}^{L}G$-局部系，可以用由全纯 ${}^{L}G$-丛 $\mathcal{E}$ 与全纯联络 $\nabla'$ 组成的二元组 $(\mathcal{E},\nabla')$ 来表示；等价地，也可以用由 $(\mathcal{E},\nabla')$ 的和乐所确定的基本群 $\pi_1(C)$ 的表示 $\rho$ 来表示。我们将主要关注\emph{不可约}的 ${}^{L}G$-局部系的情形。相应地，$\Bun_G$ 上的 $\CD$-模可以更具体地描述为偏微分方程组，其形式为本征值方程 $D_i f= E_if$，其中 $D_i$ 是 $\Bun_G$ 上一族由 Hitchin 可积系统的 Hamilton 量量子化而得的微分算符。一个更具雄心的几何 Langlands 对应版本已在 \cite{AG} 中给出，在其中有必要把对应推广到若干类可约局部系。

几何 Langlands 对应的一些原始进路，从 ${}^{L}G$-局部系为\emph{oper}$^{g)}$ 的情形出发；所谓 oper，是二元组 $(\mathcal{E},\nabla')$，其中 $\nabla'$ 在规范等价意义下可取某种标准形式。全部 oper 构成的空间，在全部局部系的模空间中形成一个 Lagrange 子空间。Beilinson 与 Drinfeld 基于共形场论的进路，对每个 oper 构造 $\Bun_G$ 上 $\CD$-模范畴中的一个对象，即仿射李代数 $\hat{\mathfrak g}_k$ 在临界水平 $k=-h^\vee$ 处的共形块。刻画共形块的 Ward 恒等式，赋予共形块层以 $\CD$-模结构。在 $k=-h^\vee$ 处，泛包络代数 $\mathcal{U}(\hat{\mathfrak g}_{k})$ 具有一个很大的中心，它同构于形式圆盘上 ${}^{L}\mathfrak{g}$-oper 的空间\cite{FF92}。由此可以证明：来自 Ward 恒等式的 $\CD$-模结构，可以用量子化 Hitchin Hamilton 量的本征值方程组 $D_i f= E_if$ 来描述，而本征值 $E_i$ 参数化 opers 的选取\cite{Fr07}。

在 \cite[Section 9.6]{Fr07} 中描述了 Beilinson--Drinfeld 几何 Langlands 对应构造的一个扩充：由 oper 的情形推广到一般的不可约局部系。其依据是这样一个事实：这类局部系总在规范等价意义下化为带有某些额外奇点的 oper\cite{Ar16}。Beilinson 与 Drinfeld 的构造对这样的 oper 给出 $\hat{\mathfrak g}_{-h^{\vee}}$ 的共形块，其中在额外奇点处插入了由 $\mathfrak{g}$ 的有限维表示诱导出的某些退化表示。在这个意义上，我们可以把一般不可约局部系的几何 Langlands 对应，看作对通常（非奇异）oper 而言已存在的对应的一种扩充。从本文将要提出的视角来看，这一观点将被证明是自然的。还应指出，在 \cite{Ga15} 为证明 \cite{AG} 中表述的强化几何 Langlands 猜想所给出的纲要中，Beilinson--Drinfeld 构造起着重要作用。

为便于后文引用，我们指出：Beilinson--Drinfeld 构造把几何 Langlands 对应中出现的 $\CD$-模描述为共形块的空间；这些空间一方面自然地纤维化于 $\Bun_G$ 之上，另一方面又纤维化于不可约 ${}^{L}G$-局部系的空间之上。

\subsection{几何 Langlands 对应——Kapustin--Witten 的进路}

Kapustin 与 Witten \cite{KW} 提出的几何 Langlands 对应版本，基于在形如 $\Sigma\times C$ 的四维流形上考察规范群为 $G_c$（$G$ 的紧实形式）的 $N=4$ SYM 理论，其中 $C$ 是 Riemann 面。在 $C$ 上紧化，可将 $N=4$ SYM 的拓扑扭曲版本有效地表示为以 $\Sigma$ 上的 Hitchin 模空间$^{(g)}$ $M_H(G)$ 为目标流形的 $2d$ 拓扑扭曲 sigma 模型。Hitchin 模空间是完全可积的，这一点在把 $M_H(G)$ 描述为环面纤维化时表现得十分明显；正是这种完全可积性，使我们能把 $N=4$ SYM 理论 S-对偶的推论，描述为联系目标流形分别为 $M_H(G)$ 与 $M_H({}^{L} G)$ 的两个 $2d$ sigma 模型的 SYZ 镜像对称。

为把上述内容与几何 Langlands 对应联系起来，Kapustin 与 Witten 考虑 $\Sigma=R\times I$、$I=[0,\pi]$ 的情形。$N=4$ SYM 理论的超对称边界条件在 $C$ 上紧化之后，将定义以 $\CM_H(G)$ 为目标流形的 $\mathcal{N}=(4,4)$ sigma 模型在 $\Sigma$ 上的边界条件。设 $\mathfrak{B}$ 为如下范畴：其对象是称为膜（brane）、保持最大数目超对称的边界条件 $\mathsf{B}$；其态射是 $\BR\times I$ 带状区域上 sigma 模型的“开弦”态空间 ${\rm Hom}_{\CM_H(G)}^{}(\mathsf{B}_1,\mathsf{B}_2)$，其中边界 $\BR\times \{0\}$ 与 $\BR\times \{\pi\}$ 分别配以边界条件 $\mathsf{B}_1$ 与 $\mathsf{B}_2$。

所谓典范余迷向膜（canonical coisotropic brane）$\mathsf{B}_{\rm cc}$ 起着特殊的作用\cite{KW,NW}。向量空间 $\mathcal{A}_{\rm cc}={\rm Hom}_{\CM_H(G)}^{}(\mathsf{B}_{\rm cc},\mathsf{B}_{\rm cc})$ 具有自然的代数结构，其乘法由“拼接开弦”定义。空间 $\CH(\mathsf{B})={\rm Hom}_{\CM_H(G)}^{}(\mathsf{B}_{\rm cc},\mathsf{B})$ 是代数 $\mathcal{A}_{\rm cc}$ 上的左模，其作用通过在带状区域 $I$ 的左边界上由 $\mathcal{A}_{\rm cc}$“拼接开弦”来定义。Kapustin 与 Witten 论证：代数 $\mathcal{A}_{\rm cc}$ 包含 $\Bun_G$ 上整体微分算符构成的代数。由此可知，空间 $\CH(\mathsf{B})$ 表示 $\Bun_G$ 上的 $\CD$-模。

把支撑在 $\BR\times\{x\}\times P$（$x\in I$，$P\in C$）上的 Wilson 与 't Hooft 线算符约化到二维 TQFT，就在膜的范畴上定义了自然的函子，它们诱导出空间 $\CH(\mathsf{B})$ 的修改。这样定义的函子在 \cite{KW} 中被等同于几何 Langlands 对应中的 Hecke 函子。对于某些膜 $\mathsf{B}$，在每个固定的 $P\in C$ 处，可将所得的修改表示为 $\CH(\mathsf{B})$ 与 ${}^{L}G$ 的某个有限维表示 $V$ 的张量积。若让点 $P\in C$ 变化所得到的一族修改能粘接成一个局部系，则称膜 $\mathsf{B}$ 满足 Hecke 本征值性质。

\cite{KW} 中指出了一族具有该性质的膜 $\mathsf{F}_\mu$。这些膜 $\mathsf{F}_\mu$ 支撑在 Hitchin 环面纤维化的纤维上。$N=4$ SYM 的 S-对偶在以 $\CM_H(G)$ 为目标流形的 sigma 模型中，表现为 SYZ 镜像对称的一个变体：它把膜 $\mathsf{F}_\mu$ 与对偶 sigma 模型（目标流形为 $\CM_H({}^{L}G)$）中的膜联系起来，后者由在 $\mu\in\CM_H({}^{L}G)$ 处具有点状支撑的摩天楼层（skyscraper sheaf）$\check{\mathsf{F}}_\mu$ 表示。

\subsection{AGT 对应——Nekrasov--Witten 的进路}

Alday、Gaiotto 与 Tachikawa 发现了某些 $N=2$ 超对称规范场论的瞬子配分函数与 Virasoro 代数共形块之间的关系\cite{AGT}。这一发现激发了大量工作，尤其引出了这类关系的各种推广。为解释 \cite{AGT} 中发现的关系，Nekrasov 与 Witten 考虑了 $\mathcal{S}$ 类的四维 $N=2$ 超对称规范场论；这类理论由极大超对称的六维量子场论在形如 $\CM^4\times C$ 的流形上沿 Riemann 面 $C$ 紧化得到。在与李代数 $\mathfrak{g}=\mathfrak{sl}_2$ 相关联的情形，所得理论具有弱耦合的 Lagrangian 描述，这些描述与 $C$ 的一种裤子分解 $\si$ 的选取相联系。对于可局部表示为纤维积 $\mathbb{R}\times I\times S^1\times S^1$ 的四维流形 $\CM^4$，\cite{NW} 论证：(i) 可以定义参数为 $\ep_1,\ep_2$ 的 $\Omega$-形变；(ii) 通过在 $S^1\times S^1$ 上紧化，可以得到用 $R\times I$ 上以 $\CM_H(G)$ 为目标流形的 sigma 模型给出的有效描述。该 sigma 模型的耦合参数为 $\ep_1/\ep_2$。

在表示 $\CM^4\simeq\mathbb{R}\times I\times S^1\times S^1$ 中，区间 $I$ 的两个端点对应于 $\CM^4$ 完全正则的点。因此，在描述 $\mathcal{S}$ 类理论在 $\CM^4$ 上紧化的、定义在 $R\times I$ 上并以 $\CM_H(G)$ 为目标流形的 sigma 模型中，必须配备特殊的边界条件。当紧化给出以 $\CM_H(G)$ 为目标流形的 sigma 模型时，\cite{NW} 论证：相应的边界条件分别由位于 $\BR\times\{0\}$ 处的典范余迷向膜的一个变体 $\mathsf{B}_{\rm cc}$，以及一类称为“oper 膜”（brane of opers）的新膜（此处记作 $\mathsf{B}_{\rm op}$）描述。\footnote{此语境下记作 $\mathsf{B}_{\rm cc}$ 的膜，与 \cite{KW} 中所考虑的膜相似但并不相同。文献 \cite{NW} 用记号 $B_{N'}$ 表示此处记作 $\mathsf{B}_{\rm op}$ 的膜。}膜 $\mathsf{B}_{\rm op}$ 是 $\mathsf{B}_{\rm cc}$ 的镜像对偶；\cite{NW} 提出，$\mathsf{B}_{\rm op}$ 是一张 Lagrange 膜，支撑在 $M_H(G)$ 内的 oper 簇上。

\cite{NW} 进一步提出，空间 $\CH={\rm Hom}_{\CM_H(G)}^{}(\mathsf{B}_{\rm cc},\mathsf{B}_{\rm op})$ 可以等同于 Virasoro 共形块的空间。为说明这一等同的动机，Nekrasov 与 Witten 指出：当 $\hbar=\ep_1/\ep_2$ 时，代数 $\mathcal{A}_{\rm cc}^{\hbar}(G)={\rm Hom}_{\CM_H(G)}^{}(\mathsf{B}_{\rm cc},\mathsf{B}_{\rm cc})$ 同构于作用在 Virasoro 共形块空间上的 Verlinde 线算符代数。镜像对称给出 $\CH(G)\simeq\CH({}^{\rm L}G)$ 的对偶描述，即空间 ${\rm Hom}_{\CM_H({}^{L}G)}^{}({\mathsf{B}}_{\rm op}',{\mathsf{B}}_{\rm cc}')$，其中 ${\mathsf{B}}_{\rm op}'$ 与 ${\mathsf{B}}_{\rm cc}'$ 分别是 ${\mathsf{B}}_{\rm op}$ 与 ${\mathsf{B}}_{\rm cc}$ 的近亲，其超对称不变性质有所改变。在对偶描述中，代数 $\check{\mathcal{A}}_{\rm cc}^{1/\hbar}({}^{L}G)={\rm Hom}_{\CM_H({}^{L}G)}^{}(\mathsf{B}_{\rm cc}',\mathsf{B}_{\rm cc}')$ 有一个显然的右作用，其作用通过在带状区域 $I$ 的右边界上“拼接开弦”来定义。$\mathcal{A}_{\rm cc}^{\hbar}(G)$ 与 $\check{\mathcal{A}}_{\rm cc}^{1/\hbar}({}^{L}G)$ 之间存在（几乎）交换的作用，这是 Virasoro 共形块空间的一个特征性质。

\subsection{反方向的联系}

Nekrasov 与 Witten 的工作\cite{NW} 与 Kapustin 与 Witten 的进路\cite{KW} 有许多共同要素，这并非偶然。一个共同的根源在于：$\CS$ 类理论与 $N=4$ SYM\cite{Wi09} 都可以看作六维 $(2,0)$ 理论在六维流形 $\CM^6=\CM^4\times C$ 上紧化的结果，其中 $C$ 是 Riemann 面，而 $\CM^4$ 是一个可局部表示为形如 $\mathbb{R}\times I \times S^1\times S^1$ 的圆纤维化的四维流形。在 $C$ 上紧化给出 $\CS$ 类理论\cite{GMN2}；在 $S^1\times S^1$ 上紧化则给出 $\mathbb{R}\times I \times C$ 上的 $N=4$ SYM，这正是 \cite{KW} 所考察的框架，\cite{Wi09} 对此作了进一步讨论。

不过应当注意，两种紧化使用的是不同的拓扑扭转，这使得结果的比较有些微妙。尽管如此，这一事实仍可用来把由 $\CS$ 类理论紧化所得的、目标流形为 $\CM_H$ 的二维 sigma 模型中的超对称边界条件，与 $C$ 上 $N=4$ SYM 中的边界条件联系起来。这些边界条件已在 Gaiotto 与 Witten 的工作\cite{GW1} 中分类。由此可以建立典范余迷向膜与 $N=4$ SYM 中纯 Neumann 边界条件之间的关系\cite{KW,GW}。交换 $S^1\times S^1$ 中的两个圆，与 $N=4$ SYM 的 S-对偶相关，而 S-对偶蕴含了 \cite{GW2} 所研究的该理论边界条件之间的关系。这使 \cite{GW} 把膜 $\mathsf{B}_{\rm op}$（即 \cite{NW} 中典范余迷向膜的镜像对偶）与 $N=4$ SYM 中由所谓 Nahm 极点边界条件下降所得的边界条件联系起来，如下文第 \ref{S:Nahm} 小节将简要讨论的。

\subsection{其他进路}

如前所述，在上面考虑的两个从六维到四维的约化中，需要使用不同的扭转。例如，要得到 \cite{NW} 所研究的框架，需要把 $\CM^4\times C$ 上的 $(2,0)$ 理论扭转成在 $\CM^4$ 上拓扑的理论；而要由六维 $(2,0)$ 理论得到 \cite{KW} 所研究的 $\BR\times I\times C$ 上拓扑扭转的 $N=4$ SYM，显然需要另一种扭转。

为描述 $\BR\times I\times C$ 上拓扑扭转的 $N=4$ SYM 的维数约化，人们自然会考虑纯拓扑的边界条件，即不依赖于 $C$ 上复结构的边界条件。这一观点促使 Ben-Zvi 与 Nadler \cite{BN16} 提出 Betti 几何 Langlands 猜想；它是 \cite{AG} 中所表述的几何 Langlands 对应的一个纯拓扑变体，自然地适用于 \cite{KW} 所研究的四维 TQFT。

Cordova 与 Jafferis 在 \cite{CJ} 中提出了理解 AGT 对应的又一个进路。它从形如 $\CM^4\times C$ 的六维流形出发，其中 $\CM^4$ 是一个压扁的四维球（squashed four-sphere），可表示为区间 $I_0$ 上的三维球纤维化，局部地 $\CM^4\sim I_0 \times S^3_{\ep_1\ep_2}$，此时 $\ep_1,\ep_2$ 是压扁参数。在 $C$ 上作部分拓扑扭转，可以论证：带有 $C$ 上扭转的 $(2,0)$ 理论的 Weyl 不变性导致配分函数既不依赖于 $\CM^4$ 的整体半径，也不依赖于 $C$ 的体积。当 $C$ 的体积很小时，有效描述将由 Pestun \cite{Pe} 及后续大量文献所研究的（压扁）四维球上 $S$ 类理论的配分函数给出。可以论证：若先沿表示 $\CM^4\sim I_0 \times S^3_{\ep_1\ep_2}$ 中出现的三维球 $S^3_{\ep_1\ep_2}$ 紧化，就能得到用 $C$ 上 CFT 表示的等价描述。以五维 SYM 理论作为中间步骤，Cordova 与 Jafferis 找到了用 $I_0\times C$ 上复 Chern--Simons 理论给出的有效描述。$C$ 两端处的边界条件由 \cite{GW1,GW2,GW} 所研究的 Nahm 极点边界条件确定。于是可以援引 Chern--Simons 理论与二维共形场论之间通常对应的推广版本，其中 Nahm 极点边界条件转化为把 WZNW 模型约化为共形 Toda 理论的约束条件。

\bigskip
\noindent\rule{\textwidth}{0.6pt}

% [Tutor Note: 易错提醒（§2.1）：oper 是二元组 $(\mathcal{E},\nabla')$，其中 $\nabla'$ 在规范等价意义下取固定的标准形；当 ${}^{L}G=SL_2$ 时，局部坐标 $z$ 下 $\nabla'=d+\begin{pmatrix}0&1\\ q&0\end{pmatrix}dz$，$q$ 为二次微分。据此可核对原文“opers 构成 Lagrange 子空间”的断言：$\dim\mathrm{Op}_{SL_2}=\dim H^0(C,K_C^2)=3g-3$，而全部 $SL_2$ 局部系的模空间维数为 $6g-6$，恰为其一半；Lagrange 性还依赖局部系模空间上的 Atiyah--Bott 辛形式。原文未给证明，完整的 oper 理论见 Drinfeld--Simpson。]
% [Tutor Note: 跳步点拨（§2.1）：$k=-h^\vee$ 是临界水平，$h^\vee$ 为对偶 Coxeter 数（$\mathfrak{sl}_2$ 时 $h^\vee=2$）。此时 Sugawara 中心荷 $c=k\dim\mathfrak g/(k+h^\vee)$ 的分母为零、Virasoro 场本身不存在，故不能按通常 CFT 方式定义共形块；Beilinson--Drinfeld 的出路正是原文所说的“很大的中心”：Feigin--Frenkel 定理给出 $\mathcal U(\hat{\mathfrak g}_{-h^\vee})$ 的中心同构于形式圆盘上 ${}^L\mathfrak g$-opers 的函数环\cite{FF92}，由此把共形块组织成 $\CD$-模，并使本征值 $E_i$ 参数化 opers\cite{Fr07}。原文 “This can be used to show” 略去的正是这一步。]
% [Tutor Note: 易错提醒（§2.2）：SYZ 镜像对称要求纤维是 special Lagrangian 环面，而不只是复环面。Hitchin 纤维化的纤维关于一个复结构 $I$ 是全纯环面，关于另一个复结构（如超 Kähler 旋转后的 $J$）才是 special Lagrangian；正是 Hitchin 模空间的超 Kähler 结构（三个复结构 $I,J,K$）保证同一纤维化可以同时充当 SYZ 纤维化。忽略超 Kähler 旋转，就会误以为“环面纤维化 + 对偶环面”已足以给出镜像对称。]
% [Tutor Note: 记号易错（§2.3）：脚注已提醒，Nekrasov--Witten 语境下的 $\mathsf{B}_{cc}$ 与 Kapustin--Witten 的 $\mathsf{B}_{cc}$ 只是“相似但不相同”，NW 原文的 $B_{N'}$ 才是本文的 $\mathsf{B}_{op}$，切勿把两篇文章的 $\mathsf{B}_{cc}$ 当成同一张膜。另注意镜像对称在交换 $G\leftrightarrow{}^LG$ 的同时把形变参数 $\hbar=\ep_1/\ep_2$ 换成 $1/\hbar$：这正是对偶代数写作 $\check{\mathcal A}_{cc}^{1/\hbar}({}^LG)$ 的原因，也是两套（几乎）交换作用能同时出现的原因。]
% [Tutor Note: 含糊点拨（§2.4）：原文说两种紧化“用了不同的拓扑扭转，使比较有些微妙”，被略去的是：六维 $(2,0)$ 理论在不同紧化中选取不同的超荷子代数作 BRST 荷——一种使 $\CM^4$ 方向拓扑化，另一种在 $\BR\times I\times C$ 上给出 Kapustin--Witten 的 GL-扭转 $N=4$ SYM——两者不能同时实现。因此比较必须借道 Gaiotto--Witten 对 $N=4$ SYM 边界条件的分类\cite{GW1,GW2}，并用 $S^1\times S^1$ 两圆交换（对应 S-对偶）把两边的边界条件对应起来；原文的“This led \cite{GW} to…”正是这一链条的终点。]
% [Tutor Note: 跳步点拨（§2.5）：原文说“Nahm 极点边界条件转化为把 WZNW 模型约化为共形 Toda 理论的约束”，这实际上就是量子 Drinfeld--Sokolov 约化：取主幂零元（在边界条件上即 Nahm 极点），把 $\hat{\mathfrak g}$-WZNW 约化为 Toda。Ben-Zvi--Nadler 的 Betti 版则是把 de Rham 侧（$\CD$-模）换成 Betti 侧（可构造层），故称“纯拓扑变体”。本节小结：BD/CFT、KW/规范场论、NW/AGT 三条进路在 §2 中只是并列陈述，真正的桥接自 §2.4 的六维共同根源与“不同扭转”开始；本文的提议（带面算符情形下 AGT 对应的推广取 Nekrasov--Shatashvili 极限，即几何 Langlands 对应）建立在 §3 及以后。]
